% Reporte de la Actividad 7
% Procesamiento secuencial y multihilo
\documentclass[12pt,oneside]{template/liverpoolthesis}
\input{template/preamble}

% ---------- Datos de esta actividad ----------
\title{Actividad 7: Procesamiento secuencial y multihilo}
\author{%
  Andrea Hernandez Huerta\\
  ID: 175523\\[1em]
  Heriberto Espino Montelongo\\
  ID: 175199\\[1em]
  \text{Profesor:} Dr. Victor Giovanni Morales Murillo}
\faculty{Universidad de las Américas Puebla (UDLAP)}
\school{Licenciatura en Ciencia de Datos}
\degree{Curso: Programación Concurrente}
\date{Fecha de entrega: 24 de septiembre de 2026}
\subject{Comparación de procesamiento secuencial y multihilo}

\addto\captionsspanish{\renewcommand{\chaptername}{Ejercicio}}
\providecommand{\figureautorefname}{Figura}
\addto\captionsspanish{\renewcommand{\figureautorefname}{Figura}}

% Inserta una fotografía proporcionada por el estudiante o deja un espacio
\newcommand{\espaciofoto}[4]{%
  \clearpage
  \begin{figure}[H]
    \centering
    \IfFileExists{figuras/#1}{%
      \includegraphics[width=0.97\textwidth,height=0.70\textheight,keepaspectratio]{figuras/#1}%
    }{%
      \fbox{%
        \begin{minipage}[c][#4][c]{0.88\textwidth}
          \centering
          \textbf{Espacio reservado para la fotografía}\\[1em]
          #2\\[1em]
          Agregar el archivo como \texttt{\detokenize{figuras/#1}}
        \end{minipage}%
      }%
    }
    \caption{#3}
  \end{figure}%
}

\begin{document}

\frontmatter
\maketitle
\tableofcontents
\listoffigures
\listoftables

\mainmatter

\chapter{Comparación secuencial y multihilo}

\textbf{Enunciado.} Comparar el rendimiento de una implementación secuencial y
una multihilo según el tamaño y la complejidad del problema. Se generan
arreglos de enteros entre 1 y 100 con 40, 100,000, 1,000,000 y 10,000,000 de
elementos. Para cada arreglo se realizan una suma sencilla y una operación de
mayor procesamiento,
\[
  \sqrt{x}\,\log(x+1),
\]
tanto de manera secuencial como con cuatro hilos. Cada implementación se
ejecuta diez veces y se registra el tiempo promedio sin incluir la generación
del arreglo.

\section{Objetivo}

El objetivo es observar cómo influyen dos variables sobre el tiempo de
ejecución: el número de elementos del arreglo y la cantidad de cómputo
realizada por elemento. La comparación permite identificar cuándo el costo de
crear, iniciar, planificar y sincronizar cuatro hilos puede ser compensado por
el trabajo que se ejecuta de forma concurrente.

\section{Entradas}

El programa utiliza cuatro tamaños de arreglo:
\[
40,\qquad 100\,000,\qquad 1\,000\,000,\qquad 10\,000\,000.
\]
Cada posición contiene un entero pseudoaleatorio entre 1 y 100. El mismo
arreglo se utiliza para comparar las dos implementaciones de una operación, de
manera que ambas reciben exactamente los mismos datos.

\section{Procesamiento}

Para la prueba de suma sencilla, la versión secuencial recorre el arreglo
completo con un solo flujo y acumula los valores en una variable
\texttt{long}. La versión multihilo divide los índices del arreglo en cuatro
intervalos. Cada hilo calcula un subtotal y lo guarda en una posición exclusiva
de un arreglo de resultados parciales.

Para la prueba de mayor procesamiento se utiliza, para cada elemento \(x\),
\[
f(x)=\sqrt{x}\log(x+1).
\]
En Java se implementa mediante
\texttt{Math.sqrt(valor) * Math.log(valor + 1.0)}. La versión secuencial
acumula todos los valores en un único flujo; la versión multihilo aplica la
misma división en cuatro intervalos y calcula cuatro subtotales.

Los cuatro hilos se crean primero y se inician mediante \texttt{start()}.
Posteriormente, el hilo principal ejecuta \texttt{join()} sobre cada uno.
Solo después de que los cuatro han terminado se combinan sus resultados
parciales.

\section{Estructuras de control utilizadas}

Se emplean ciclos \texttt{for} para generar los datos, recorrer los arreglos y
repetir cada medición diez veces. La versión multihilo utiliza un arreglo de
cuatro objetos \texttt{Thread}. Cada hilo recibe una tarea mediante una
expresión lambda y procesa el intervalo
\[
\left[
\frac{i n}{4},
\frac{(i+1)n}{4}
\right),
\]
donde \(i\in\{0,1,2,3\}\) es el identificador del hilo y \(n\) es el tamaño del
arreglo.

Cada hilo escribe únicamente en su propia posición del arreglo de resultados
parciales. Por ello no se requiere un cerrojo para los subtotales. Las llamadas
a \texttt{join()} garantizan que todos los resultados parciales estén
disponibles antes de combinarlos.

\section{Medición del tiempo}

La generación del arreglo se realiza antes de iniciar cualquier cronómetro, tal
como solicita la actividad. Cada implementación se ejecuta diez veces. Para
cada repetición se utiliza \texttt{System.nanoTime()} y se acumula la duración.
El promedio en milisegundos se calcula como
\[
\bar t =
\frac{\sum_{k=1}^{10} t_k}{10\cdot 10^6}.
\]

El tiempo de la versión multihilo incluye la creación, el inicio y la
sincronización de los hilos, pues esas operaciones forman parte de la
implementación que se desea comparar.

\section{Comprobación de resultados}

Para la suma sencilla, los resultados secuencial y multihilo deben coincidir
exactamente. Para la operación con \texttt{double}, el orden de las sumas
cambia entre ambas implementaciones y puede generar diferencias mínimas de
redondeo. Por ello se comprueba equivalencia con una tolerancia relativa de
\(10^{-10}\). Si alguna comprobación falla, el programa lanza una excepción en
lugar de presentar tiempos de implementaciones que producen resultados
distintos.

\section{Resultados}

Los tiempos son dependientes del procesador, de la JVM y de la carga del
sistema. La siguiente tabla debe completarse con los promedios obtenidos al
ejecutar el programa en la computadora utilizada para la evidencia.

\begin{table}[H]
\centering
\small
\caption{Comparación de tiempos promedio entre las implementaciones.}
\label{tab:resultados}
\begin{tabular}{r l c c c}
\hline
\textbf{Tamaño} & \textbf{Operación} &
\textbf{Secuencial (ms)} & \textbf{Multihilo (ms)} & \textbf{Más rápido}\\
\hline
40 & Suma & --- & --- & ---\\
40 & Mayor procesamiento & --- & --- & ---\\
100,000 & Suma & --- & --- & ---\\
100,000 & Mayor procesamiento & --- & --- & ---\\
1,000,000 & Suma & --- & --- & ---\\
1,000,000 & Mayor procesamiento & --- & --- & ---\\
10,000,000 & Suma & --- & --- & ---\\
10,000,000 & Mayor procesamiento & --- & --- & ---\\
\hline
\end{tabular}
\end{table}

\espaciofoto{codigo.png}
{Fotografía del código Java de la Actividad 7}
{Código de la comparación secuencial y multihilo.}
{9cm}

\espaciofoto{salida.png}
{Fotografía de la ejecución completa y de la tabla de tiempos obtenida en la computadora del estudiante}
{Salida real del programa con los tiempos promedio.}
{9cm}

\chapter{Análisis y conclusiones}

\section{¿En qué casos fue más rápida cada implementación?}

La respuesta definitiva debe basarse en los valores medidos de la
\autoref{tab:resultados}. Para arreglos pequeños, el trabajo útil puede ser tan
breve que el costo de crear, iniciar y sincronizar cuatro hilos sea comparable
o incluso mayor que el propio cálculo. En ese caso es razonable observar un
menor tiempo en la implementación secuencial.

Al crecer el arreglo, especialmente en la operación
\(\sqrt{x}\log(x+1)\), existe más cómputo independiente que repartir entre los
cuatro hilos. Si la computadora dispone de capacidad de ejecución paralela,
la versión multihilo puede entonces reducir el tiempo promedio. Esta
interpretación debe contrastarse con los datos obtenidos en la ejecución
local.

\section{¿Cómo influyeron el tamaño del arreglo y la complejidad de la operación?}

Aumentar el tamaño incrementa el número de elementos que deben procesarse.
Además, la suma sencilla realiza una operación aritmética muy económica,
mientras que la segunda prueba evalúa una raíz cuadrada y un logaritmo para
cada elemento. Por tanto, el segundo caso asigna más trabajo computacional a
cada sección del arreglo.

Cuando el trabajo por elemento es pequeño, los costos asociados con los hilos
representan una fracción importante del tiempo total. Al incrementar el tamaño
o la complejidad, esa fracción puede disminuir porque cada hilo permanece más
tiempo realizando trabajo útil.

\section{¿Por qué utilizar varios hilos no siempre mejora el rendimiento?}

Crear hilos y coordinarlos tiene un costo. Además, el sistema operativo debe
planificar su ejecución y el procesador dispone de una cantidad finita de
núcleos y recursos compartidos. Si la tarea es demasiado pequeña, el trabajo
adicional de administrar los hilos puede superar el tiempo que se pretende
ahorrar. Por ello, el uso de varios hilos no implica por sí solo una mejora de
rendimiento; el beneficio depende del tamaño del problema, de la complejidad
del cálculo y del hardware donde se ejecuta.

\section{Conclusión}

La actividad compara dos formas de ejecutar exactamente el mismo cálculo. La
versión secuencial evita el costo de coordinación, mientras que la versión
multihilo divide el trabajo en cuatro secciones independientes y sincroniza su
terminación mediante \texttt{join()}. La comparación correcta requiere medir
ambas alternativas sobre los mismos datos, excluir la generación del arreglo
del cronómetro y verificar que los resultados numéricos sean equivalentes.

Los datos de la \autoref{tab:resultados} permiten determinar experimentalmente
en qué punto el trabajo disponible es suficiente para compensar el costo de
usar cuatro hilos.

\cleardoublepage
\phantomsection
\addcontentsline{toc}{chapter}{Referencias}
\nocite{*}
\bibliography{template/references}

\end{document}
